\label{sec:formula}

\subsection{Zone Assignment}

For each tract $t \in \mathcal{T}$ and each zone type $z \in \{1, \ldots, 6\}$,
we define $\mathrm{zone}_z(t)$ as the identifier of the zone of type $z$ that
contains the centroid of tract $t$.
This centroid-in-polygon assignment is the standard approach for assigning
tracts to administrative zones; it fails only when a zone boundary bisects a
tract so narrowly that the centroid falls on the wrong side.
For practical purposes, this case is rare: census tracts are designed to respect
administrative boundaries wherever possible, and in most states school district
and county subdivision boundaries are co-extensive with or contained within
tract boundaries.

Let $Z$ be the set of zone types for which data is available in the state being
processed.
The minimum $|Z| = 3$ is achieved when only school district, county subdivision,
and electric utility data are available; the maximum $|Z| = 6$ is achieved when
fire district, water district, and police precinct data are also present.

\subsection{Zone Co-membership Score}

For two adjacent tracts $u, v \in \mathcal{T}$, the zone co-membership score is
\begin{equation}
  \mathrm{zone\_score}(u,v)
  \;=\;
  \frac{|\{z \in Z : \mathrm{zone}_z(u) = \mathrm{zone}_z(v)\}|}{|Z|}
  \;\in\; [0,\, 1].
  \label{eq:zone-score}
\end{equation}
The numerator counts how many of the available zone types are shared between $u$
and $v$; the denominator normalizes to the number of available zone types.
The score is 1 when all available zones are shared, 0 when no zones are shared.

\paragraph{Interpretation.}
A score of 1 means that tracts $u$ and $v$ send their children to the same schools,
share a fire station, buy water from the same utility, purchase electricity from the
same provider, and are governed by the same county subdivision---they are members
of the same administrative community in every dimension available.
A score of 0 means that $u$ and $v$ share no administrative zone; they may be
adjacent geographically but are administratively in completely different communities.
A score of $1/3$ (when $|Z| = 3$) might arise when two tracts share school district
and county subdivision but are served by different electric utilities---a common
situation at the edges of utility service territories in rural areas.

\subsection{Edge Weight Modifier}

The zone co-membership score is applied as an \emph{additive boost}:
\begin{equation}
  w(u,v)
  \;=\;
  w_{\mathrm{base}}(u,v) \;\times\; \bigl(1 + \alpha \cdot \mathrm{zone\_score}(u,v)\bigr),
  \quad \alpha = 1.0 \text{ (default)},
  \label{eq:zone-weight}
\end{equation}
where $w_{\mathrm{base}}(u,v)$ is the base edge weight from the active base mode.

\begin{remark}[Relationship to M.0 multiplicative formula]
The M.0 framework defines $w = w_{\mathrm{base}} \times \mathrm{sim}$, which \emph{reduces}
weights for dissimilar pairs (minimum: $0$ when $\mathrm{sim}=0$).
M.6's additive boost instead \emph{increases} weights for co-zone pairs
(minimum: $w_{\mathrm{base}}$ when $\mathrm{zone\_score}=0$; maximum: $2 \times w_{\mathrm{base}}$
when $\mathrm{zone\_score}=1$ and $\alpha=1$).
Both are valid instantiations of the M-track community character framework:
multiplicative (M.1--M.5, M.7) models similarity as an attenuation of existing weights;
additive boost (M.6) models co-membership as a strengthening signal on top of
existing weights. For zone co-membership, the additive interpretation is natural:
two tracts in the same fire district have a \emph{stronger} connection than baseline,
rather than a less-attenuated one.
\end{remark}

\paragraph{Boundary cases.}
When $\mathrm{zone\_score}(u,v) = 1$ (all zones shared) and $\alpha = 1$:
\begin{equation}
  w(u,v) = w_{\mathrm{base}}(u,v) \times 2 \quad \text{(maximum bond)}.
\end{equation}
When $\mathrm{zone\_score}(u,v) = 0$ (no zones shared):
\begin{equation}
  w(u,v) = w_{\mathrm{base}}(u,v) \times 1 = w_{\mathrm{base}}(u,v) \quad \text{(baseline, unmodified)}.
\end{equation}
This design choice is deliberate: tracts that share no administrative zone type
are not penalized relative to the base weight.
The algorithm incentivizes preserving administratively co-membered pairs without
discriminating against pairs that happen to lie at administrative boundaries.

\paragraph{Alpha parameter.}
The $\alpha$ parameter controls the strength of the zone co-membership incentive.
At $\alpha = 1$ (default), all-zones-shared tracts have twice the base cut cost.
At $\alpha = 2$, all-zones-shared tracts have three times the base cut cost.
At $\alpha = 0$, the formula reduces to $w(u,v) = w_{\mathrm{base}}(u,v)$ with no
zone effect.
In practice, $\alpha = 1$ is the recommended value because it provides a meaningful
but not dominating incentive: METIS can still make a cut between co-membered tracts
if the geometric cost of avoiding the cut is too high.

\subsection{Mathematical Properties}

\paragraph{Symmetry.}
$\mathrm{zone\_score}(u,v) = \mathrm{zone\_score}(v,u)$ because set equality is
symmetric: $\mathrm{zone}_z(u) = \mathrm{zone}_z(v)$ iff $\mathrm{zone}_z(v) = \mathrm{zone}_z(u)$.
Consequently $w(u,v) = w(v,u)$, satisfying METIS's undirected graph requirement.

\paragraph{Non-negativity.}
$\mathrm{zone\_score}(u,v) \geq 0$ always (it is a ratio of non-negative integers),
so $w(u,v) \geq w_{\mathrm{base}}(u,v) \geq 0$ for all $\alpha \geq 0$.

\paragraph{Boundedness.}
$\mathrm{zone\_score}(u,v) \leq 1$, so $w(u,v) \leq (1 + \alpha) \times w_{\mathrm{base}}(u,v)$.
For $\alpha = 1$, this means no edge weight is more than doubled.

\paragraph{Graceful degradation.}
If zone data is entirely unavailable ($|Z| = 0$), the formula is undefined.
The implementation guards against this by always including county co-membership
as a fallback dimension when no other zone data is available.
County co-membership can always be computed from Census tract FIPS codes (the
first five digits of the tract FIPS are the state and county FIPS).
In the extreme case where no special district data is available and county subdivision
is also unavailable, county co-membership serves as $Z = \{$county$\}$ with $|Z| = 1$.

\subsection{Implementation Invariants}

The following invariants are enforced as L0 unit tests:
\begin{itemize}
  \item \texttt{all\_zones\_shared\_doubles\_weight}:
    $\mathrm{zone\_score} = 1.0 \Rightarrow w = 2 \times w_{\mathrm{base}}$ when $\alpha = 1.0$.
  \item \texttt{no\_zones\_shared\_baseline}:
    $\mathrm{zone\_score} = 0.0 \Rightarrow w = w_{\mathrm{base}}$ exactly.
  \item \texttt{zone\_score\_range\_0\_to\_1}:
    $\mathrm{zone\_score}(u,v) \in [0.0, 1.0]$ for all adjacent tract pairs.
  \item \texttt{school\_district\_split\_never\_increases}:
    For any state tested, the number of school district boundary crossings under
    zone-membership weights is $\leq$ the number under geographic weights.
  \item \texttt{zone\_score\_symmetric}:
    $\mathrm{zone\_score}(u,v) = \mathrm{zone\_score}(v,u)$ for all pairs.
\end{itemize}

\subsection{CLI Activation}

The zone co-membership weight mode is activated with:
\begin{verbatim}
bisect state --state NC --year 2020 \
  --weights-override zone-membership \
  --version v_zone_m6
\end{verbatim}
The alpha parameter can be tuned via the plan YAML configuration:
\begin{verbatim}
algorithm:
  weights: zone-membership
  zone_alpha: 1.0
  zone_types: [school_district, county_subdivision, electric_utility,
               fire_district, water_district]
\end{verbatim}
